Having established in Section~\ref{sec:2} a common mixture representation for conditional redshift distributions and a corresponding family of estimators, we now specify the experimental design choices required to implement them on realistic datasets. These choices define the conditions under which the calibration performance of the estimators is assessed in Section~\ref{sec:4} across the ensemble of realised experiments described below. In that section, however, the conditional redshift distributions are still evaluated one experiment at a time. Section~\ref{sec:5} then changes perspective by treating the ensemble itself as the object over which statistical and systematic uncertainty is propagated and marginalised.

Our implementation builds on the SOM-based augmentation pipeline developed in Paper~I, but extends its role from point-estimate photometric redshifts to the calibration of conditional redshift distributions \(\phi_m(z)\). To keep the present paper readable as a stand-alone description of the full analysis pipeline, Sections~\ref{sec:3.1} and~\ref{sec:3.2} briefly summarise the main ingredients inherited from Paper~I, while referring the reader there for detailed implementation choices and validation tests. In particular, Section~\ref{sec:3.1} summarises the catalogue construction across the experimental ensemble, and Section~\ref{sec:3.2} describes the photometric-redshift modelling and the definitions of the lens and source galaxy populations. We then introduce the additional ingredients required for the present focus on calibrating tomographic redshift distributions: Section~\ref{sec:3.3} defines the tomographic \texttt{Target} and \texttt{Reference} samples, together with the common SOM-based discretisation and hierarchical clustering used for conditional summarisation, while Section~\ref{sec:3.4} specifies the calibration configurations considered in this work.

\subsection{Dataset construction across experiments} \label{sec:3.1}

Following Paper~I, we organise the data pipeline into four catalogue types: the \texttt{Application} datasets (the photometric population used for cosmological analyses), the spectroscopic \texttt{Degradation} datasets (observational analogues with selection effects), the simulated \texttt{Augmentation} datasets (used to fill gaps in sparsely sampled regions of feature space), and the merged \texttt{Combination} datasets. For each LSST depth scenario, Year~1 (Y1) and Year~10 (Y10), we generate a baseline experiment (indexed by \(i=0\)) together with \(I = 500\) additional statistical variations (indexed by \(i = 1, \ldots, I\)) that incorporate randomised photometric noise, footprint variations, and non-uniform spectroscopic selection functions, yielding \(I+1 = 501\) realisations per survey scenario. In Paper~I, this experimental ensemble was introduced and validated in the context of point-estimate photometric-redshift calibration; here we adopt the same ensemble as the experimental backbone for evaluating the conditional redshift distributions \(\phi_m(z)\) and, later, for constructing the experiment-marginalised uncertainty framework of Section~\ref{sec:5}.

The \texttt{Application} datasets are designed to reproduce LSST-like photometric samples based on the \texttt{OpenUniverse2024} catalogue \citep{2025arXiv250105632O}, incorporating realistic noise characteristics and flux-uncertainty models \citep{2024AJ....168...80C}. Following Paper~I, all galaxies are required to satisfy a combined \(r{+}i\)-band signal-to-noise ratio \(\mathrm{SNR}_{i+r} > 10\), to eliminate very faint extended sources, and a resolution cut \(\bigl(R^{i+r}/R_\mathrm{PSF}^{i+r}\bigr)^{2} > 0.1\), where \(R^{i+r}\) and \(R_\mathrm{PSF}^{i+r}\) are the galaxy and PSF sizes measured from the coadded \(r{+}i\) images, to reduce contamination from unresolved point sources. In addition, we exclude objects brighter than \(16 \) in \(i-\)band to mimic the masking of bright, saturated sources, and remove galaxies fainter than \(24.05\) for Y1 and \(25.30\) for Y10, in accordance with the LSST DESC SRD. 

The \texttt{Degradation} datasets originate from the same parent population but are assembled to mimic compilations of multiple heterogeneous spectroscopic surveys. Each survey component is specified by magnitude, colour, and redshift selection criteria, together with additional sampling that reflects the corresponding spectroscopic success rates. The resulting spectroscopic compilation is therefore heterogeneous by construction, capturing the diversity of spectroscopic coverage and incompleteness anticipated from present and upcoming observing programmes.

To mitigate incompleteness in the spectroscopic feature-space coverage, we generate synthetic galaxies from the independent \texttt{CosmoDC2} catalogue \citep{2019ApJS..245...26K}. These galaxies are mapped into the photometric feature space of the \texttt{Application} datasets using a pre-trained SOM (hereafter the \textit{augmentation SOM}) implemented with \texttt{Somoclu} \citep{2013arXiv1305.1422W}. In our default augmentation set-up, the augmentation SOM has resolution \(140 \times 140\) for Y1 and \(180 \times 180\) for Y10, and is trained on the baseline \texttt{Application} datasets to capture the target photometric manifold. To reduce differences in key photometric properties between the \texttt{CosmoDC2} and \texttt{OpenUniverse2024} samples, we perform SOM-based weighted resampling to align their distributions in feature space, followed by adaptive, targeted filling to supplement underpopulated or empty SOM cells. These steps yield the \texttt{Augmentation} datasets and follow the procedure developed in Paper~I; here we retain only the summary needed to define the datasets that enter the calibration of tomographic redshift distributions.

The two catalogues are deliberately built on distinct stellar-population modelling prescriptions. \texttt{OpenUniverse2024} populates galaxies with the auto-differentiable \texttt{Diffsky} framework, assembling continuous parametric models for halo mass assembly \citep{2021OJAp....4E...7H}, star formation histories \citep{2023MNRAS.518..562A}, and stellar-population synthesis \citep{2023MNRAS.521.1741H}. \texttt{CosmoDC2}, by contrast, couples the \texttt{Galacticus} semi-analytic model \citep{2012NewA...17..175B} to the \texttt{UniverseMachine} empirical halo--galaxy connection \citep{2019MNRAS.488.3143B} through \texttt{GalSampler} \citep{2020MNRAS.495.5040H}, drawing stellar populations from a more restricted array of discrete templates. As demonstrated in Paper~I, these differing prescriptions yield visibly different colour--redshift manifolds, particularly at faint magnitudes and high redshifts. This contrast is precisely what motivates our setup: using \texttt{OpenUniverse2024} as the photometric proxy and \texttt{CosmoDC2} as the augmentation source makes the residual mismatch between the two simulations a controlled, tractable proxy for the irreducible mismatch between any augmentation catalogue and real LSST observations.

We emphasise that even in a real LSST analysis, the augmentation catalogue would necessarily be drawn from simulated data, since augmentation is designed precisely to fill regions of colour space where spectroscopic coverage is unavailable. The mismatch between simulated and observed galaxy populations therefore remains a fundamental limitation of any augmentation strategy and is not an artefact of our validation set-up; we return to this point in Section~\ref{sec:6.1}.

The \texttt{Combination} datasets are constructed by merging the \texttt{Degradation} samples with the \texttt{Augmentation} populations, providing more complete coverage of the galaxy SED manifold and an alternative training set for photometric-redshift modelling under the augmentation framework. In Paper~I, we showed that this strategy reduces systematic biases in point-estimate redshifts, particularly for high-redshift galaxies. In the present work, we use the same catalogue infrastructure as the starting point for calibrating conditional redshift distributions \(\phi_m(z)\).

For controlled stress-testing purposes, we additionally construct the \texttt{Restriction} datasets by resampling \texttt{CosmoDC2} galaxies to match the \texttt{Application} SOM distribution and then downsampling to the same sample size as the corresponding \texttt{Combination} datasets. This yields a calibration sample of matched size but based entirely on simulated galaxies, enabling a direct assessment of calibration performance when the colour--redshift relation is intentionally mismatched to the target photometric population (Section~\ref{sec:3.4}). Table~\ref{tab:1} summarises the sample types, their origin and construction, roles in the pipeline, and implementation status in a real LSST analysis.

\begin{table*}
    \centering
    \begin{tabular}{p{2.4cm}p{4.2cm}p{5.8cm}p{3.2cm}}
        \toprule
        Sample & Origin and construction & Role in pipeline & LSST implementation \\
        \midrule
        \texttt{Application} & \texttt{OpenUniverse2024}; photometric selection & Photometric target population for cosmological analysis & Replaced by observed photometry \\
        \texttt{Degradation} & \texttt{OpenUniverse2024}; subsampled by magnitude, colour, redshift, and success-rate criteria & Spectroscopic calibration compilations with heterogeneous selection effects & Replaced by spectroscopic compilations \\
        \texttt{Augmentation} & \texttt{CosmoDC2}; SOM-weighted resampling with adaptive cell filling & Simulated galaxies supplementing sparse regions of feature space & Simulated by construction \\
        \texttt{Combination} & \texttt{Degradation} + \texttt{Augmentation}; merged & Combined calibration sample with improved feature-space coverage & Merged observed and simulated set \\
        \texttt{Restriction} & \texttt{CosmoDC2}; resampled to match the \texttt{Application} SOM distribution and downsampled to \texttt{Combination} size & Controlled sample for stress-testing under simulation--observation mismatch & Validation only \\
        \texttt{Target} & \texttt{Application}; tomographic bin selection & Photometric galaxies for which \(\phi_m(z)\) is calibrated & Observed, tomographically selected \\
        \texttt{Reference} & Configuration dependent; see Table~\ref{tab:2} & Calibration sample with secure redshifts for cluster-level conditional density estimation & Configuration dependent \\
        \bottomrule
    \end{tabular}
    \caption{Summary of the sample types used in this work. Each row specifies the parent catalogue and construction procedure, the role within the calibration pipeline, and the expected implementation in a real LSST analysis. Samples labelled ``Replaced by'' would be substituted with observational data; ``Simulated by construction'' indicates that augmentation necessarily relies on simulated catalogues even in an operational setting (Section~\ref{sec:3.1}); ``Validation only'' denotes samples retained solely for controlled diagnostic tests.}
\label{tab:1}
\end{table*}

\subsection{Photometric redshift estimation and definition of galaxy populations} \label{sec:3.2}

Photometric redshifts are estimated using the \texttt{FlexZBoost} algorithm \citep{2020A&C....3000362D}, which models the conditional redshift probability density function (PDF) directly from broadband photometry. The method represents each conditional PDF as a linear combination of Fourier basis functions, with random-forest regression employed to predict the corresponding coefficients. This formulation enables flexible and non-parametric estimation of the conditional redshift densities for individual galaxies \citep{10.1214/17-EJS1302}, while maintaining robustness against complex colour--redshift degeneracies.

Hyperparameter optimisation and model training are carried out within the Redshift Assessment Infrastructure Layers (\texttt{RAIL}\footnote{\url{https://github.com/LSSTDESC/rail}}) framework \citep{2025arXiv250502928T}. The trained models are subsequently applied to the \texttt{Application} datasets to generate posterior redshift distributions for individual galaxies, thereby ensuring consistency between the training-data construction and the inference stage. In the present work, our emphasis is on how these posteriors enter the calibration of conditional redshift distributions and the associated sample definitions. The specific training datasets used for \texttt{FlexZBoost} are configuration dependent and specified in Section~\ref{sec:3.4}.

For the purposes of \(3\times2\)pt cosmological analyses, galaxies are divided into foreground lens and background source populations. Source galaxies are further selected following \citet{2013MNRAS.434.2121C} to maximise the fidelity of weak gravitational-lensing measurements, requiring per-galaxy shape measurement uncertainty smaller than the intrinsic ellipticity dispersion (\(\sigma_\gamma < \sigma_\epsilon\)) and photometric redshifts in the range \(0.05 < z_{\text{phot}} < 2.95\). Lens galaxies are defined as a brighter foreground subsample used for clustering and galaxy--galaxy lensing analyses. Following \citet{2021PhRvD.103d3503P}, a redshift-dependent magnitude cut \(i < 4 \cdot z_{\text{phot}} + 18\) is imposed within the range \(0.2 < z_{\text{phot}} < 1.2\), ensuring sufficiently precise photometric-redshift estimates for this population. The selection cuts on \(r\)+\(i\) imaging quality (Section~\ref{sec:3.1}) and on per-galaxy shape uncertainty mirror real-survey practice for constructing shape catalogues, while photometric-redshift estimation, calibration SOM training, and conditional density summarisation use the full LSST six-band photometry (\(ugrizy\)) for both lens and source samples.

\subsection{Tomographic calibration samples and infrastructure for conditional summarisation} \label{sec:3.3}

In this subsection, we define the tomographic calibration samples and the SOM-based infrastructure used to construct conditional redshift summaries within each bin. Throughout, tomographic bin assignments are defined using maximum-a-posteriori (MAP) photometric-redshift point estimates,
\begin{equation} \label{eq:20}
    z_{\text{phot}} \equiv \arg\max_{z}\, p_{\mathrm{ML}}(z \mid \mathbf{x}),
\end{equation}
where \(p_{\mathrm{ML}}(z \mid \mathbf{x})\) is the \texttt{FlexZBoost} posterior. For source galaxies, we partition the sample into five equal-number tomographic bins in each of the 501 experiments under both the Y1 and Y10 survey scenarios. For lens galaxies, we adopt equal-width redshift bins over the range \([0.2,\,1.2]\): five bins for Y1 and ten bins for Y10. The finer lens binning in Y10 enables a more detailed characterisation of the redshift evolution of large-scale structure. This binning strategy follows the provisional benchmark of the LSST~DESC SRD, although refinements---including optimised source-bin assignments explored in the LSST~DESC Tomography Optimization Challenge \citep{2021OJAp....4E..13Z} and neural-network-based lens-binning improvements \citep{2023ApJ...950...49M}---may be explored in future work.

These tomographic assignments define the bin-dependent selection functions \(s_m\). Within each tomographic bin, we define the photometric galaxies drawn from the \texttt{Application} datasets---either lenses or sources---as the \texttt{Target} samples for which conditional redshift distributions are summarised. The \texttt{Reference} samples are the calibration sets used to infer the colour--redshift relation from secure redshifts. Depending on the configuration, the \texttt{Reference} samples are drawn from the spectroscopic \texttt{Degradation} datasets, from the augmented \texttt{Combination} datasets, or, in stress tests, from the fully simulated \texttt{Restriction} datasets. These configuration-dependent choices are described in Section~\ref{sec:3.4}.

Following the estimator definitions in Section~\ref{sec:2.2}, \texttt{Reference} samples may be constructed either tomographically or globally. A common practical choice is to construct tomographic \texttt{Reference} samples for the lens population by applying the lens-galaxy selection \(s_m\) in each bin, reflecting the relatively narrow redshift range and brighter selection of this sample. For the source population, the dominant challenge is to capture a global colour--redshift relation across a broader and fainter SED manifold. One may therefore construct a global source \texttt{Reference} sample using all spectroscopic galaxies that satisfy the source selection criteria, independent of tomographic binning. This reduces the susceptibility to bin-to-bin fluctuations in the spectroscopic subset induced by imperfect point-estimate assignments. Alternative definitions, including swaps between tomographic and global constructions for lenses and sources, are specified in Section~\ref{sec:3.4}.

To discretise photometric feature space, we employ a second SOM (hereafter the \textit{calibration SOM}) and compute Best Matching Units (BMUs) for all galaxies in the \texttt{Target} and \texttt{Reference} samples. The dataset used to train the calibration SOM, together with its resolution, need not coincide with the augmentation SOM of Section~\ref{sec:3.1}. These choices are treated as analysis-design parameters and are specified for each configuration in Section~\ref{sec:3.4}. After training, we reduce the high-resolution calibration SOM using agglomerative clustering \citep{2011arXiv1109.2378M}, merging neighbouring cells into larger clusters containing galaxies with similar SEDs. Hereafter, we adopt a number of clusters equal to one quarter of the total SOM cells for lens galaxies and one half for source galaxies, reflecting the smaller size of the bright lens population. This coarsening increases the number of galaxies per cluster and mitigates stochastic noise from sparsely occupied cells.

Throughout this work, we rely exclusively on LSST-like observations and do not incorporate additional information from the LSST Deep Drilling Fields, such as deeper multi-band photometry or near-infrared coverage. Consequently, a DES-style multi-layer SOM architecture \citep[e.g.][]{2019MNRAS.489..820B, 2021MNRAS.505.4249M, 2024arXiv240800922C, 2025arXiv250907964G, 2025arXiv251023566Y} is not directly applicable in our context and is deferred to future study. In this respect, our calibration strategy is closer in spirit to KiDS-style approaches \citep[e.g.][]{2020A&A...637A.100W, 2025arXiv250319440W}, which anchor colour--redshift relations using spectroscopic information within a single-SOM framework.

Given the resulting cluster assignment \(c \in \{1, \ldots, C\}\), the \texttt{Target} and \texttt{Reference} samples can be partitioned into the same set of clusters. The cluster-level conditional densities \(H_c(z_n \mid s_m)\) are then constructed using the estimators defined in Section~\ref{sec:2.2}, and combined with the mixture weights \(\Pi_{m,c}\) to obtain \(\phi_m(z_n)\) via Equation~(\ref{eq:10}). Clusters that contain \texttt{Target} galaxies but no \texttt{Reference} counterparts are excluded from the mixture by assigning zero weight \(\Pi_{m,c}=0\), ensuring that the calibration is supported only in regions of feature space anchored by the calibration sample. After applying this support criterion, we renormalise the remaining \(\Pi_{m,c}\) so that \(\sum_c \Pi_{m,c}=1\) on the supported set. This criterion is applied uniformly across the \texttt{DIR} and \texttt{Stack} estimators, enabling a like-for-like comparison and a well-defined \texttt{Hybrid} combination. The same criterion is also imposed when constructing the \texttt{Truth} benchmark, so that all summaries are defined on an identical supported set of clusters.

\subsection{Varying configurations for calibration assessment} \label{sec:3.4}

We now define a suite of calibration configurations that differ both in the datasets adopted and in the roles assigned to them in: (i) training the photometric-redshift point estimates and conditional PDFs, (ii) training the calibration SOM used for discretisation and clustering in feature space, and (iii) constructing the \texttt{Reference} samples used to estimate the cluster-level conditional densities \(H_c(z_n \mid s_m)\). These configurations represent alternative analysis choices rather than a ranking of performance. Across all cases, conditional redshift-distribution summaries are constructed using the common procedure described in Sections~\ref{sec:3.1}--\ref{sec:3.3}: dataset construction, \texttt{FlexZBoost} modelling, lens and source sample definitions with tomographic binning, BMU assignment on a trained calibration SOM, agglomerative clustering to define the cluster index \(c\), mixture construction, and the cluster-support criterion.

At a conceptual level, \texttt{Gold}, \texttt{Silver}, and \texttt{Titanium} trace a progression from a purely spectroscopic calibration pipeline to a fully augmented fiducial framework. The \texttt{Copper} and \texttt{Iron} configurations then isolate, respectively, the sensitivity to the dataset used to train the calibration SOM and to the construction of the \texttt{Reference} samples. Finally, \texttt{Zinc} provides an explicit stress test in which the colour--redshift relation is intentionally mismatched to the target photometric population. The six configurations and their dataset choices are summarised in Table~\ref{tab:2}.

\begin{table*}
    \centering
    \begin{tabular}{p{4.5cm}cccccc}
        \toprule
        Configuration name & \texttt{Gold} & \texttt{Silver} & \texttt{Titanium} & \texttt{Copper} & \texttt{Iron} & \texttt{Zinc} \\
        \midrule
        Dataset for training calibration SOM & \texttt{Degradation} & \texttt{Degradation} & \texttt{Application} & \texttt{Combination} & \texttt{Application} & \texttt{Restriction} \\
        Dataset for training photometric redshifts & \texttt{Degradation} & \texttt{Combination} & \texttt{Combination} & \texttt{Combination} & \texttt{Combination} & \texttt{Restriction} \\
        Dataset for defining \texttt{Reference} samples & \texttt{Degradation} & \texttt{Degradation} & \texttt{Combination} & \texttt{Combination} & \texttt{Combination} & \texttt{Restriction} \\
        Definition of \texttt{Reference} samples of lens galaxies & tomographic & tomographic & tomographic & tomographic & global & tomographic \\
        Definition of \texttt{Reference} samples of source galaxies & global & global & global & global & tomographic & global \\
        Applicability & Realistic & Realistic & Fiducial & Realistic & Realistic & Stress test \\
        \bottomrule
    \end{tabular}
    \caption{Summary of the six calibration configurations considered in this work. The configurations differ in the datasets used to train the calibration SOM and photometric-redshift models, and in the construction of the \texttt{Reference} samples for lens and source galaxies. The bottom row indicates the applicability of each configuration: the first five (\texttt{Gold}--\texttt{Iron}) represent physically motivated analysis choices applicable to real LSST data, with \texttt{Titanium} serving as the fiducial augmentation-enabled set-up, while \texttt{Zinc} is an intentional stress test using deliberately mismatched simulated training data (Section~\ref{sec:3.4}).}
\label{tab:2}
\end{table*}

\subsubsection{\texttt{Gold}: configuration without data augmentation} \label{sec:3.4.1}

The \texttt{Gold} configuration excludes simulated galaxies at all stages. Photometric redshifts are trained using only the spectroscopic \texttt{Degradation} datasets, from which \texttt{FlexZBoost} produces both MAP point estimates and conditional PDFs. These MAP point estimates define the tomographic bins following the procedure in Section~\ref{sec:3.3}.

For discretisation, the calibration SOM is trained on the \texttt{Degradation} datasets, adopting a reduced resolution relative to the default augmentation SOM to reflect the smaller number of spectroscopically selected galaxies (\(100 \times 100\) for Y1 and \(125 \times 125\) for Y10). We compute BMUs for galaxies in both the \texttt{Degradation} and \texttt{Application} datasets, so that \texttt{Target} samples are partitioned in the same feature space as the spectroscopic \texttt{Reference} set. Cluster-level conditional densities are then estimated using \texttt{Degradation} as the \texttt{Reference} sample: \texttt{DIR} summarises secure spectroscopic redshifts within each cluster, while \texttt{Stack} aggregates the conditional PDFs trained on \texttt{Degradation}; these are combined through the \texttt{Hybrid} construction.

\subsubsection{\texttt{Silver}: augmentation only for point estimates} \label{sec:3.4.2}

The \texttt{Silver} configuration isolates the impact of augmented galaxies on tomographic binning while retaining a purely spectroscopic calibration reference for the \texttt{DIR} component. Relative to \texttt{Gold}, photometric-redshift point estimates and conditional PDFs are trained on the augmented \texttt{Combination} datasets, and the resulting MAP point estimates are used to redefine tomographic bins (Section~\ref{sec:3.3}).

All steps that rely on secure redshifts follow \texttt{Gold}: the calibration SOM is trained on \texttt{Degradation} at the same reduced resolution, BMUs are computed for \texttt{Degradation} and \texttt{Application}, and the \texttt{Reference} samples used for \texttt{DIR} are drawn from \texttt{Degradation}. The \texttt{Stack} component, however, is constructed from the conditional PDFs trained on \texttt{Combination}, ensuring consistency with the point-estimate assignments. In this way, \texttt{Silver} separates the effect of augmentation on sample definition from its effect on the calibration relation entering \texttt{DIR}.

\subsubsection{\texttt{Titanium}: data augmentation along all stages} \label{sec:3.4.3}

The \texttt{Titanium} configuration provides the fiducial configuration of our SOM-based augmentation framework, in which simulated galaxies are incorporated consistently in both photometric-redshift training and the construction of the calibration \texttt{Reference} samples (Table~\ref{tab:2}). Photometric-redshift point estimates and conditional PDFs are trained on the \texttt{Combination} datasets, and the resulting MAP point estimates define the tomographic bins as in Section~\ref{sec:3.3}.

We adopt the default augmentation SOM of Section~\ref{sec:3.1} (resolution \(140 \times 140\) for Y1 and \(180 \times 180\) for Y10) as the calibration SOM as well, yielding a single feature-space partition that both defines the augmentation manifold and provides the cluster assignments used for conditional summarisation. The \texttt{Reference} samples are drawn from \texttt{Combination} and are constructed tomographically for lenses and globally for sources (Table~\ref{tab:2}), providing support for the \texttt{DIR} estimator. The same \texttt{Combination}-trained \texttt{FlexZBoost} posteriors are stacked within clusters to form the \texttt{Stack} estimator, ensuring that augmentation enters the summarisation consistently at all stages.

\subsubsection{\texttt{Copper}: calibration SOM trained on \texttt{Combination} datasets} \label{sec:3.4.4}

The \texttt{Copper} configuration tests the sensitivity of the discretisation step to the dataset used to train the calibration SOM under the augmentation framework. Relative to \texttt{Titanium}, the only change is that the calibration SOM is trained on the \texttt{Combination} datasets (Table~\ref{tab:2}), while photometric-redshift point estimates and conditional PDFs remain trained on \texttt{Combination}. BMUs are then computed for all \texttt{Target} and \texttt{Reference} samples using this retrained calibration SOM. All other choices follow \texttt{Titanium}: \texttt{Reference} samples are drawn from \texttt{Combination}, constructed tomographically for lenses and globally for sources, and the same mixture construction and support criterion are applied. Conceptually, \texttt{Titanium} trains the calibration SOM on \texttt{Application}, so that the feature-space partition is anchored by the photometric target population, whereas \texttt{Copper} retrains it on \texttt{Combination}, so that the partition is instead defined on the augmented feature distribution used by the calibration sample. This contrast isolates how much the choice of the SOM-training data---as opposed to the \texttt{Reference}-sample construction, which is varied separately by \texttt{Iron}---affects the resulting calibration.

\subsubsection{\texttt{Iron}: swapping \texttt{Reference} sample construction} \label{sec:3.4.5}

The \texttt{Iron} configuration isolates the impact of constructing \texttt{Reference} samples globally versus tomographically. Relative to \texttt{Titanium}, we swap the \texttt{Reference} definitions by adopting a global \texttt{Reference} sample for lenses and tomographic \texttt{Reference} samples for sources (Table~\ref{tab:2}). All other choices follow \texttt{Titanium}: photometric redshifts are trained on \texttt{Combination}, the discretisation uses the same default calibration SOM (identical to the augmentation SOM, as in \texttt{Titanium}) and clustering procedure, and the mixture construction and support criterion are unchanged. This configuration therefore targets the sensitivity of \(H_c(z_n\mid s_m)\) to the \texttt{Reference} construction while holding fixed both the feature-space partition and the point-estimate binning.

\subsubsection{\texttt{Zinc}: stress test with fully simulated training data} \label{sec:3.4.6}

The \texttt{Zinc} configuration provides an extreme stress test in which the training data for both the photometric-redshift estimators and the calibration SOM are drawn from the \texttt{Restriction} datasets. These catalogues are based on fully simulated galaxies (e.g.\ \texttt{CosmoDC2}) and therefore need not share the same colour--redshift relation as the \texttt{Application} datasets (e.g.\ \texttt{OpenUniverse2024}). The \texttt{Reference} samples are also constructed from \texttt{Restriction} (Table~\ref{tab:2}).

All remaining choices follow \texttt{Titanium}, including the lens/source selections, the \texttt{Reference} definitions (tomographic for lenses and global for sources), and the cluster-support criterion described in Section~\ref{sec:3.3}. This configuration quantifies the degradation expected when the training and calibration relations are mismatched to the target photometric population.